\chapter{习题选解}

% ===================== 第一节 =====================
\section{第一章习题（群论）}

\begin{example}[若群 $G$ 的每个元都适合 $x^2=e$，则 $G$ 是交换群]
若群 $G$ 的每一个元都适合方程 $x^2=e$，那么 $G$ 是交换群。

\textbf{解：} 令 $a$ 和 $b$ 是 $G$ 的任意两个元。由题设
\[
(ab)(ab)=(ab)^2=e
\]
另一方面
\[
(ab)(ba)=ab^2a=aea=a^2=e
\]
于是有 $(ab)(ab)=(ab)(ba)$。利用消去律，得
\[
ab=ba
\]
所以 $G$ 是交换群。
\end{example}

\begin{example}[有限群里阶大于 $2$ 的元个数为偶数]
在一个有限群里阶大于 $2$ 的元的个数一定是偶数。

\textbf{解：} 设 $G$ 是有限群，$a\in G$ 的阶大于 $2$。则 $a^{-1}\ne a$（否则 $a^2=e$ 意味着 $|a|\le2$），故 $a^{-1}$ 也是阶大于 $2$ 的元。因此阶大于 $2$ 的元总是成对出现，个数为偶数。
\end{example}

\begin{example}[阶为偶数的有限群中阶等于 $2$ 的元个数为奇数]
假定 $G$ 是一个阶为偶数的有限群。在 $G$ 里阶等于 $2$ 的元的个数一定是奇数。

\textbf{解：} 由上题，阶大于 $2$ 的元成对出现，共有偶数个。$G$ 的元分为三类：单位元 $e$（$1$ 个），阶等于 $2$ 的元（设 $k$ 个），阶大于 $2$ 的元（偶数个）。由于 $|G|$ 是偶数，$1+k+\text{偶数}$ 为偶数，故 $k$ 是奇数。
\end{example}

\begin{example}[有限群的元的阶有限]
一个有限群的每一个元的阶都有限。

\textbf{解：} 设 $G$ 是有限群，$a\in G$。考察序列 $a,a^2,a^3,\ldots$，由于 $G$ 有限，必有 $i<j$ 使 $a^i=a^j$，即 $a^{j-i}=e$。故 $|a|\le j-i<\infty$。
\end{example}

\begin{example}[$|a|=|a^{-1}|$]
设 $G$ 是群，$a\in G$，证明 $a$ 与 $a^{-1}$ 的阶相同。

\textbf{解：} 设 $|a|=n$，则 $(a^{-1})^n=(a^n)^{-1}=e^{-1}=e$，故 $|a^{-1}|\le n$。设 $|a^{-1}|=m$，则 $a^m=((a^{-1})^m)^{-1}=e^{-1}=e$，故 $|a|\le m$。因此 $m=n$，即 $|a|=|a^{-1}|$。
\end{example}

\begin{example}[模 $12$ 的剩余类加群的所有子群]
找出模 $12$ 的剩余类加群 $G=\mathbb{Z}_{12}$ 的所有子群。

\textbf{解：} $G$ 是阶为 $12$ 的循环群，其子群都是循环群。容易看出：
\begin{align*}
([0])&=\{[0]\}\\
([1])=([5])=([7])=([11])&=G\\
([2])=([10])&=\{[2],[4],[6],[8],[10],[0]\}\\
([3])=([9])&=\{[3],[6],[9],[0]\}\\
([4])=([8])&=\{[4],[8],[0]\}\\
([6])&=\{[6],[0]\}
\end{align*}
以上是 $G$ 的所有子群。
\end{example}

% ===================== 第二节 =====================
\section{第二章习题（环论）}

\begin{example}[环的中心是交换子环]
证明：一个环的中心是一个交换子环。

\textbf{解：} 设 $R$ 是环，$Z(R)=\{z\in R:zr=rz,\,\forall r\in R\}$ 是 $R$ 的中心。
\begin{enumerate}[(i)]
  \item $0\in Z(R)$；
  \item 若 $z_1,z_2\in Z(R)$，则 $(z_1-z_2)r=z_1r-z_2r=rz_1-rz_2=r(z_1-z_2)$，故 $z_1-z_2\in Z(R)$；
  \item 若 $z_1,z_2\in Z(R)$，则 $(z_1z_2)r=z_1(rz_2)=(rz_1)z_2=r(z_1z_2)$，故 $z_1z_2\in Z(R)$；
  \item $Z(R)$ 中任意两元可交换（由定义直接得）。
\end{enumerate}
故 $Z(R)$ 是 $R$ 的一个交换子环。
\end{example}

\begin{example}[除环的中心是域]
证明：一个除环的中心是一个域。

\textbf{解：} 设 $D$ 是除环，$Z(D)$ 是 $D$ 的中心。由上题 $Z(D)$ 是交换子环。若 $z\in Z(D)$，$z\ne0$，则 $z$ 在 $D$ 中有逆元 $z^{-1}$。对任意 $r\in D$，由 $zr=rz$ 两边乘以 $z^{-1}$ 得 $rz^{-1}=z^{-1}r$，故 $z^{-1}\in Z(D)$。因此 $Z(D)$ 是域。
\end{example}

\begin{example}[$\mathbb{Q}$ 是 $\mathbb{Q}(i)$ 的唯一真子域]
证明有理数域 $\mathbb{Q}$ 是所有复数 $a+bi$（$a,b$ 是有理数）作成的域 $\mathbb{Q}(i)$ 的唯一的真子域。

\textbf{解：} 显然 $\mathbb{Q}$ 是 $\mathbb{Q}(i)$ 的一个真子域。设 $F$ 是 $\mathbb{Q}(i)$ 的任一子域，那么 $F$ 含有元 $a\ne0$，因而含有 $a^{-1}a=1$，由此 $F$ 含有一切整数和一切有理数，所以 $\mathbb{Q}\subset F\subset\mathbb{Q}(i)$。若 $F\ne\mathbb{Q}$，那么至少有一个数 $a+bi\in F$，$a,b\in\mathbb{Q}$，$b\ne0$。于是 $a+bi-a=bi\in F$，$b^{-1}\cdot bi=i\in F$。由此 $F$ 含有一切 $a+bi$ 而 $F=\mathbb{Q}(i)$。所以 $\mathbb{Q}$ 是 $\mathbb{Q}(i)$ 的唯一的真子域。
\end{example}

\begin{example}[$(4)$ 是 $\mathbb{Z}$ 的极大理想且 $\mathbb{Z}/(4)$ 有零因子]
设 $R=\mathbb{Z}$，证明主理想 $(4)$ 是 $R$ 的极大理想，并说明 $R/(4)$ 有零因子从而不是一个域。

\textbf{解：} $(4)$ 刚好含有一切 $4n$（$n\in\mathbb{Z}$）。设 $\mathfrak{U}$ 是 $R$ 的一个理想，$(4)\subsetneq\mathfrak{U}$。那么 $\mathfrak{U}\ni a=4q+2\nmid4\mathbb{Z}$，由此 $\mathfrak{U}\ni a-4q=2$，而 $\mathfrak{U}\ni2m$ 对所有 $m\in\mathbb{Z}$。故 $\mathfrak{U}=(2)=R$，说明 $(4)$ 是极大理想。

在 $R/(4)$ 中，$[2]\ne[0]$ 而 $[2][2]=[4]=[0]$。因此 $R/(4)$ 有零因子而不是域。
\end{example}

% ===================== 第三节 =====================
\section{第三章习题（整环分解）}

\begin{example}[求 $\mathbb{Z}[i]$ 中 $5$ 的因子]
求 $\mathbb{Z}[i]$ 中 $5$ 的因子。

\textbf{解：} 设 $a+bi\mid5$，则 $\exists\,c+di\in\mathbb{Z}[i]$ 使 $(a+bi)(c+di)=5$，故
\[
25=(a^2+b^2)(c^2+d^2),
\]
即 $a^2+b^2\in\{1,5,25\}$。

$5$ 的\textbf{平凡因子}（单位及其相伴元）：$1,\,-1,\,i,\,-i,\;5,\,-5,\,5i,\,-5i$。

$5$ 的\textbf{真因子}（对应 $a^2+b^2=5$）：$\pm2\pm i,\;\pm1\pm2i$。
\end{example}

\begin{example}[$\mathbb{Z}[i]/(1+i)$ 的元素个数]
设 $R=\mathbb{Z}[i]$（高斯整数环），求 $R/(1+i)$ 有多少个元。

\textbf{解：} 先看主理想 $(1+i)$ 含有哪些元。设 $a+bi\in(1+i)$，那么
\[
a+bi=(x+yi)(1+i)=(x-y)+(x+y)i
\]
因此若 $x,y$ 同奇同偶，则 $a,b$ 同为偶数；若 $x,y$ 一奇一偶，则 $a,b$ 同为奇数。总之，$a+bi\in(1+i)$ 当且仅当 $a,b$ 有相同的奇偶性。

因此 $R/(1+i)$ 中代表元只需从 $a,b$ 奇偶性不同的高斯整数中取，共有两个陪集 $[0]$（$a,b$ 同奇偶）和 $[1]$（$a,b$ 异奇偶）。故 $R/(1+i)$ 共有 $\mathbf{2}$ 个元。
\end{example}
